\documentclass[11pt,reqno]{amsart}
\usepackage{amsmath}
\usepackage{amssymb}
\usepackage{amsthm}
\usepackage{enumerate}
\usepackage[notcite,notref]{showkeys}
\usepackage[usenames]{color}
\usepackage{eepic,epic}
\usepackage{url}
\usepackage{kotex}
\usepackage{epstopdf}

\usepackage{graphicx}

\usepackage{hyperref}

% PAGE SETTING ---------------------------------------

\textheight 22.5  true cm
\textwidth 15 true cm
\voffset -1.0 true cm
\hoffset -1.0 true cm
\marginparwidth= 2 true cm
\renewcommand{\baselinestretch}{1.2}

% THEOREMS -------------------------------------------

\renewcommand{\thefootnote}{\fnsymbol{footnote}}
\newtheorem{thm}{Theorem}[section]
\newtheorem{cor}[thm]{Corollary}
\newtheorem{lem}[thm]{Lemma}
\newtheorem{prop}[thm]{Proposition}
\newtheorem{slem}[thm]{Sublemma}
\theoremstyle{definition}
\newtheorem{defn}[thm]{Definition}
\newtheorem{step}{Step}
\theoremstyle{remark}
\newtheorem{rem}[thm]{Remark}
\newtheorem{prob}[thm]{Problem}
\numberwithin{equation}{section}

\newtheorem{innercustomthm}{Theorem}
\newenvironment{customthm}[1]
{\renewcommand\theinnercustomthm{#1}\innercustomthm}
{\endinnercustomthm}


\newcommand{\ep}{\epsilon}
\newcommand{\lx}{{\lambda, \xi}}
\newcommand{\pa}{\partial}
\renewcommand{\(}{\left(}
\renewcommand{\)}{\right)}

% MATH ------------------------------------------------

\newcommand{\R}{\mathbb{R}}
\newcommand{\inner}[2]{\left< #1, #2 \right>}
\newcommand{\test}[1]{\left[ #1 \right]}
\newcommand{\setdefine}{ \ \big| \ }

%-----------------------------------------------------------------

\newenvironment{text-box}{}{}
\newcommand{\tags}[1]{#1}
\newcommand{\subheading}[1]{#1}

% ----------------------------------------------------------------
\begin{document}
\title[]
%{Convergence rate of numerical scheme finding extremal functions of Sobolev inequalities}
{test - 3}

\date{2025-04-16}

\maketitle

\tags{
    supremum, infimum, bounded, approximation property, 
    completeness axiom, archimedean principle, density of rationals, 
    reflection principle, monotone property
}

\section{test}

\subsection{test - 3}

This is test. Note that $S^n := \{ x \in \R^{n+1} \setdefine \| x \| = 1 \}$.

\subheading{Definition}

\begin{equation}
    3 \epsilon + 2 \lx - \pa f(x) = \test{Hello}
\end{equation}

\begin{align}
    \( \R \in \inner{f}{g} \) \\
    3 \ep + 2 \lx - \pa f(x) = \test{Hello}
\end{align}

\begin{text-box}
    1) This is test. Note that $S^n := \{ x \in \R^{n+1} \setdefine \| x \| = 1 \}$.

    2) This is test. Note that $S^n := \{ x \in \R^{n+1} \setdefine \| x \| = 1 \}$. \\
    3) This is test. Note that $S^n := \{ x \in \R^{n+1} \setdefine \| x \| = 1 \}$.

    \begin{equation}
        ax + by = c
    \end{equation}
    
    4) This is test. Note that $S^n := \{ x \in \R^{n+1} \setdefine \| x \| = 1 \}$.
\end{text-box}

\end{document}

